
\documentclass{article}
\usepackage{colortbl}
\usepackage{makecell}
\usepackage{multirow}
\usepackage{supertabular}

\begin{document}

\newcounter{utterance}

\centering \large Interaction Transcript for game `clean\_up', experiment `1\_medium\_3obj\_en', episode 0 with qwen3{-}next{-}80b{-}a3b{-}instruct{-}t1.0.
\vspace{24pt}

{ \footnotesize  \setcounter{utterance}{1}
\setlength{\tabcolsep}{0pt}
\begin{supertabular}{c@{$\;$}|p{.15\linewidth}@{}p{.15\linewidth}p{.15\linewidth}p{.15\linewidth}p{.15\linewidth}p{.15\linewidth}}
    \# & $\;$A & \multicolumn{4}{c}{Game Master} & $\;\:$B\\
    \hline

    \theutterance \stepcounter{utterance}  
    & & \multicolumn{4}{p{0.6\linewidth}}{
        \cellcolor[rgb]{0.9,0.9,0.9}{
            \makecell[{{p{\linewidth}}}]{
                \texttt{\tiny{[P1$\langle$GM]}}
                \texttt{Sono il tuo game master, e stai giocando un gioco collaborativo con la seguente griglia come tavolo da gioco:} \\
\\ 
\texttt{```} \\
\texttt{ 1234567} \\
\texttt{╔═╤═╤═╤═╗} \\
\texttt{║◌│◌│◌│◌║ 1} \\
\texttt{║◌└─┼─┴─╢ 2} \\
\texttt{║◌◌◌│◌◌◌║ 3} \\
\texttt{╟──┬┴───╢ 4} \\
\texttt{║◌◌│P◌◌◌║ 5} \\
\texttt{║LC│◌◌◌◌║ 6} \\
\texttt{║◌◌│◌◌◌◌║ 7} \\
\texttt{╚══╧════╝} \\
\\ 
\texttt{```} \\
\\ 
\texttt{* Il bordo superiore mostra le coordinate x che aumentano verso destra, e il bordo destro le coordinate y che aumentano verso il basso.} \\
\texttt{* I seguenti oggetti sono posizionati casualmente sulla tua griglia: 'C', 'L', 'P'.} \\
\\ 
\\ 
\texttt{L'altro giocatore vede una variante del tavolo da gioco, dove gli oggetti sono posizionati in diverse posizioni casuali. Non puoi vedere il tavolo dell'altro giocatore, e loro non possono vedere il tuo.} \\
\\ 
\\ 
\texttt{**Obiettivo:**} \\
\\ 
\texttt{Entrambi i giocatori devono spostare gli oggetti sui rispettivi sfondi in modo che oggetti identici finiscano alle stesse coordinate. Devi comunicare con l'altro giocatore per concordare una configurazione obiettivo comune.} \\
\\ 
\\ 
\texttt{**Regole:**} \\
\\ 
\texttt{* In ogni turno, puoi inviare esattamente uno dei seguenti due comandi:} \\
\texttt{ 1. `COMUNICA: <MESSAGE>`: per inviare un messaggio (tutto fino alla prossima interruzione di riga) all'altro giocatore. Lo inoltrerò al tuo partner.} \\
\texttt{2. `SPOSTA: <OBJECT>, (<X>, <Y>)`: per spostare un oggetto in una nuova posizione, dove <X> è la colonna e <Y> è la riga. Ti farò sapere se la tua mossa è andata a buon fine o no.} \\
\texttt{* Se non rispetti il formato, o invii diversi comandi contemporaneamente, devo penalizzarti.} \\
\texttt{* Se entrambi i giocatori accumulano più di 8 penalità, perdete entrambi il gioco.} \\
\texttt{* È vitale che comunichi con l'altro giocatore riguardo al tuo stato obiettivo! L'*unico* modo in cui puoi trasmettere la tua strategia all'altro giocatore è usando il comando `COMUNICA: <MESSAGE>`!} \\
\\ 
\\ 
\texttt{**Spostamento Oggetti**} \\
\\ 
\texttt{* Puoi spostare oggetti solo in celle entro i limiti della griglia. La cella di destinazione deve essere vuota, cioè deve contenere solo il simbolo '◌'.} \\
\texttt{* Se provi a spostare un oggetto in un posto che non è vuoto, o provi a spostarlo fuori dalla griglia, devo penalizzarti. Hai un altro tentativo.} \\
\texttt{* Prima di fare una mossa, controlla due volte che il posto di destinazione sia vuoto, e non contenga alcuna lettera, cornice, o linea!} \\
\\ 
\\ 
\texttt{**Fine del Gioco**} \\
\\ 
\texttt{Se pensi di aver raggiunto l'obiettivo di allineare tutti gli oggetti, puoi chiedere all'altro giocatore di finire il gioco inviando COMUNICA: finito?. Se l'altro giocatore ti chiede di finire il gioco, e rispondi con COMUNICA: finito!, the game will end.} \\
\\ 
\texttt{Entrambi i giocatori vincono se il gioco finisce entro 12 round, dove un round è definito come due giocatori che inviano ciascuno un comando valido.} \\
\\ 
\texttt{**Punteggio:**} \\
\\ 
\texttt{Più vicini sono gli oggetti identici in entrambi i tavoli da gioco, più punti ottieni. Le penalità riducono i tuoi punti. Riesci a battere il record?} \\
\\ 
\\ 
\texttt{Per favore invia un messaggio all'altro giocatore per iniziare il gioco!} \\
            }
        }
    }
    & \\ \\

    \theutterance \stepcounter{utterance}  
    & \multicolumn{4}{p{0.6\linewidth}}{
        \cellcolor[rgb]{0.9,0.9,0.9}{
            \makecell[{{p{\linewidth}}}]{
                \texttt{\tiny{[P1$\rangle$GM]}}
                \texttt{COMUNICA: Ciao! Il mio tavolo ha 'P' in (5,5), 'L' in (2,6), e 'C' in (1,6). Ho una griglia 7x7. Per allinearci, suggerisco di portare tutti gli oggetti verso il centro: P in (4,4), L in (3,4), C in (5,4). Che ne pensi? Possiamo scambiarci le posizioni precise dei nostri oggetti e coordinarci. Attendo tuo feedback!} \\
            }
        }
    }
    & & \\ \\

    \theutterance \stepcounter{utterance}  
    & & & \multicolumn{2}{p{0.3\linewidth}}{
        \cellcolor[rgb]{0.9,0.9,0.9}{
            \makecell[{{p{\linewidth}}}]{
                \texttt{\tiny{[GM$|$GM]}}
                \texttt{COMUNICA: Ciao! Il mio tavolo ha 'P' in (5,5), 'L' in (2,6), e 'C' in (1,6). Ho una griglia 7x7. Per allinearci, suggerisco di portare tutti gli oggetti verso il centro: P in (4,4), L in (3,4), C in (5,4). Che ne pensi? Possiamo scambiarci le posizioni precise dei nostri oggetti e coordinarci. Attendo tuo feedback!} \\
            }
        }
    }
    & & \\ \\

    \theutterance \stepcounter{utterance}  
    & & \multicolumn{4}{p{0.6\linewidth}}{
        \cellcolor[rgb]{0.9,0.9,0.9}{
            \makecell[{{p{\linewidth}}}]{
                \texttt{\tiny{[GM$\rangle$P2]}}
                \texttt{Sono il tuo game master, e stai giocando un gioco collaborativo con la seguente griglia come tavolo da gioco:} \\
\\ 
\texttt{```} \\
\texttt{ 1234567} \\
\texttt{╔═╤═╤═╤═╗} \\
\texttt{║◌│◌│L│◌║ 1} \\
\texttt{║◌└─┼─┴─╢ 2} \\
\texttt{║◌◌◌│◌◌C║ 3} \\
\texttt{╟──┬┴───╢ 4} \\
\texttt{║◌◌│◌◌◌◌║ 5} \\
\texttt{║◌◌│◌P◌◌║ 6} \\
\texttt{║◌◌│◌◌◌◌║ 7} \\
\texttt{╚══╧════╝} \\
\\ 
\texttt{```} \\
\\ 
\texttt{* Il bordo superiore mostra le coordinate x che aumentano verso destra, e il bordo destro le coordinate y che aumentano verso il basso.} \\
\texttt{* I seguenti oggetti sono posizionati casualmente sulla tua griglia: 'C', 'L', 'P'.} \\
\\ 
\\ 
\texttt{L'altro giocatore vede una variante del tavolo da gioco, dove gli oggetti sono posizionati in diverse posizioni casuali. Non puoi vedere il tavolo dell'altro giocatore, e loro non possono vedere il tuo.} \\
\\ 
\\ 
\texttt{**Obiettivo:**} \\
\\ 
\texttt{Entrambi i giocatori devono spostare gli oggetti sui rispettivi sfondi in modo che oggetti identici finiscano alle stesse coordinate. Devi comunicare con l'altro giocatore per concordare una configurazione obiettivo comune.} \\
\\ 
\\ 
\texttt{**Regole:**} \\
\\ 
\texttt{* In ogni turno, puoi inviare esattamente uno dei seguenti due comandi:} \\
\texttt{ 1. `COMUNICA: <MESSAGE>`: per inviare un messaggio (tutto fino alla prossima interruzione di riga) all'altro giocatore. Lo inoltrerò al tuo partner.} \\
\texttt{2. `SPOSTA: <OBJECT>, (<X>, <Y>)`: per spostare un oggetto in una nuova posizione, dove <X> è la colonna e <Y> è la riga. Ti farò sapere se la tua mossa è andata a buon fine o no.} \\
\texttt{* Se non rispetti il formato, o invii diversi comandi contemporaneamente, devo penalizzarti.} \\
\texttt{* Se entrambi i giocatori accumulano più di 8 penalità, perdete entrambi il gioco.} \\
\texttt{* È vitale che comunichi con l'altro giocatore riguardo al tuo stato obiettivo! L'*unico* modo in cui puoi trasmettere la tua strategia all'altro giocatore è usando il comando `COMUNICA: <MESSAGE>`!} \\
\\ 
\\ 
\texttt{**Spostamento Oggetti**} \\
\\ 
\texttt{* Puoi spostare oggetti solo in celle entro i limiti della griglia. La cella di destinazione deve essere vuota, cioè deve contenere solo il simbolo '◌'.} \\
\texttt{* Se provi a spostare un oggetto in un posto che non è vuoto, o provi a spostarlo fuori dalla griglia, devo penalizzarti. Hai un altro tentativo.} \\
\texttt{* Prima di fare una mossa, controlla due volte che il posto di destinazione sia vuoto, e non contenga alcuna lettera, cornice, o linea!} \\
\\ 
\\ 
\texttt{**Fine del Gioco**} \\
\\ 
\texttt{Se pensi di aver raggiunto l'obiettivo di allineare tutti gli oggetti, puoi chiedere all'altro giocatore di finire il gioco inviando COMUNICA: finito?. Se l'altro giocatore ti chiede di finire il gioco, e rispondi con COMUNICA: finito!, the game will end.} \\
\\ 
\texttt{Entrambi i giocatori vincono se il gioco finisce entro 12 round, dove un round è definito come due giocatori che inviano ciascuno un comando valido.} \\
\\ 
\texttt{**Punteggio:**} \\
\\ 
\texttt{Più vicini sono gli oggetti identici in entrambi i tavoli da gioco, più punti ottieni. Le penalità riducono i tuoi punti. Riesci a battere il record?} \\
\\ 
\texttt{L'altro giocatore ha iniziato il gioco inviando questo messaggio:} \\
\\ 
\texttt{"Ciao! Il mio tavolo ha 'P' in (5,5), 'L' in (2,6), e 'C' in (1,6). Ho una griglia 7x7. Per allinearci, suggerisco di portare tutti gli oggetti verso il centro: P in (4,4), L in (3,4), C in (5,4). Che ne pensi? Possiamo scambiarci le posizioni precise dei nostri oggetti e coordinarci. Attendo tuo feedback!"} \\
\\ 
\texttt{Qual è il tuo primo comando?} \\
            }
        }
    }
    & \\ \\

    \theutterance \stepcounter{utterance}  
    & & & \multicolumn{4}{p{0.6\linewidth}}{
        \cellcolor[rgb]{0.9,0.9,0.9}{
            \makecell[{{p{\linewidth}}}]{
                \texttt{\tiny{[GM$\langle$P2]}}
                \texttt{COMUNICA: Ciao! Il mio tavolo ha 'P' in (5,6), 'L' in (3,1), e 'C' in (7,3). La tua proposta di allinearli al centro (P in (4,4), L in (3,4), C in (5,4)) è ottima! Concordo pienamente. Procedo subito a spostare i miei oggetti verso quei target. Inizio con L: lo muovo da (3,1) a (3,4). Verificherò che (3,4) sia vuota prima di spostare.} \\
            }
        }
    }
     \\ \\

    \theutterance \stepcounter{utterance}  
    & & & \multicolumn{2}{p{0.3\linewidth}}{
        \cellcolor[rgb]{0.9,0.9,0.9}{
            \makecell[{{p{\linewidth}}}]{
                \texttt{\tiny{[GM$|$GM]}}
                \texttt{COMUNICA: Ciao! Il mio tavolo ha 'P' in (5,6), 'L' in (3,1), e 'C' in (7,3). La tua proposta di allinearli al centro (P in (4,4), L in (3,4), C in (5,4)) è ottima! Concordo pienamente. Procedo subito a spostare i miei oggetti verso quei target. Inizio con L: lo muovo da (3,1) a (3,4). Verificherò che (3,4) sia vuota prima di spostare.} \\
            }
        }
    }
    & & \\ \\

    \theutterance \stepcounter{utterance}  
    & & \multicolumn{4}{p{0.6\linewidth}}{
        \cellcolor[rgb]{0.9,0.9,0.9}{
            \makecell[{{p{\linewidth}}}]{
                \texttt{\tiny{[P1$\langle$GM]}}
                \texttt{Il tuo messaggio è stato inoltrato all'altro giocatore.} \\
\\ 
\\ 
\texttt{Stai attualmente giocando il round 1 di massimo 12.} \\
\texttt{Avete accumulato collettivamente 0 di 8 penalità.} \\
\texttt{L'altro giocatore ha inviato questo messaggio:} \\
\\ 
\\ 
\texttt{"Ciao! Il mio tavolo ha 'P' in (5,6), 'L' in (3,1), e 'C' in (7,3). La tua proposta di allinearli al centro (P in (4,4), L in (3,4), C in (5,4)) è ottima! Concordo pienamente. Procedo subito a spostare i miei oggetti verso quei target. Inizio con L: lo muovo da (3,1) a (3,4). Verificherò che (3,4) sia vuota prima di spostare."} \\
\\ 
\\ 
\texttt{Qual è il tuo prossimo comando?} \\
            }
        }
    }
    & \\ \\

    \theutterance \stepcounter{utterance}  
    & \multicolumn{4}{p{0.6\linewidth}}{
        \cellcolor[rgb]{0.9,0.9,0.9}{
            \makecell[{{p{\linewidth}}}]{
                \texttt{\tiny{[P1$\rangle$GM]}}
                \texttt{SPOSTA: L, (3,6)} \\
            }
        }
    }
    & & \\ \\

    \theutterance \stepcounter{utterance}  
    & & & \multicolumn{2}{p{0.3\linewidth}}{
        \cellcolor[rgb]{0.9,0.9,0.9}{
            \makecell[{{p{\linewidth}}}]{
                \texttt{\tiny{[GM$|$GM]}}
                \texttt{SPOSTA: L, (3,6)} \\
            }
        }
    }
    & & \\ \\

    \theutterance \stepcounter{utterance}  
    & & & \multicolumn{2}{p{0.3\linewidth}}{
        \cellcolor[rgb]{0.9,0.9,0.9}{
            \makecell[{{p{\linewidth}}}]{
                \texttt{\tiny{[GM$|$GM]}}
                \texttt{Penalità: (3,6) non è vuoto, ma contiene '│'.} \\
\\ 
\texttt{Avete accumulato collettivamente 1 di 8 penalità.} \\
\texttt{Riprova!} \\
            }
        }
    }
    & & \\ \\

    \theutterance \stepcounter{utterance}  
    & & \multicolumn{4}{p{0.6\linewidth}}{
        \cellcolor[rgb]{0.9,0.9,0.9}{
            \makecell[{{p{\linewidth}}}]{
                \texttt{\tiny{[P1$\langle$GM]}}
                \texttt{Penalità: (3,6) non è vuoto, ma contiene '│'.} \\
\\ 
\texttt{Avete accumulato collettivamente 1 di 8 penalità.} \\
\texttt{Riprova!} \\
            }
        }
    }
    & \\ \\

    \theutterance \stepcounter{utterance}  
    & \multicolumn{4}{p{0.6\linewidth}}{
        \cellcolor[rgb]{0.9,0.9,0.9}{
            \makecell[{{p{\linewidth}}}]{
                \texttt{\tiny{[P1$\rangle$GM]}}
                \texttt{SPOSTA: L, (2,5)} \\
            }
        }
    }
    & & \\ \\

    \theutterance \stepcounter{utterance}  
    & & & \multicolumn{2}{p{0.3\linewidth}}{
        \cellcolor[rgb]{0.9,0.9,0.9}{
            \makecell[{{p{\linewidth}}}]{
                \texttt{\tiny{[GM$|$GM]}}
                \texttt{SPOSTA: L, (2,5)} \\
            }
        }
    }
    & & \\ \\

    \theutterance \stepcounter{utterance}  
    & & & \multicolumn{2}{p{0.3\linewidth}}{
        \cellcolor[rgb]{0.9,0.9,0.9}{
            \makecell[{{p{\linewidth}}}]{
                \texttt{\tiny{[GM$|$GM]}}
                \texttt{Spostato 'L' in (2,5) con successo. La tua griglia aggiornata appare così:} \\
\texttt{```} \\
\texttt{ 1234567} \\
\texttt{╔═╤═╤═╤═╗} \\
\texttt{║◌│◌│◌│◌║ 1} \\
\texttt{║◌└─┼─┴─╢ 2} \\
\texttt{║◌◌◌│◌◌◌║ 3} \\
\texttt{╟──┬┴───╢ 4} \\
\texttt{║◌L│P◌◌◌║ 5} \\
\texttt{║◌C│◌◌◌◌║ 6} \\
\texttt{║◌◌│◌◌◌◌║ 7} \\
\texttt{╚══╧════╝} \\
\texttt{```} \\
            }
        }
    }
    & & \\ \\

    \theutterance \stepcounter{utterance}  
    & & \multicolumn{4}{p{0.6\linewidth}}{
        \cellcolor[rgb]{0.9,0.9,0.9}{
            \makecell[{{p{\linewidth}}}]{
                \texttt{\tiny{[GM$\rangle$P2]}}
                \texttt{Il tuo messaggio è stato inoltrato all'altro giocatore.} \\
\\ 
\\ 
\texttt{Stai attualmente giocando il round 2 di massimo 12.} \\
\texttt{Avete accumulato collettivamente 1 di 8 penalità.} \\
\texttt{L'altro giocatore ha spostato un oggetto sulla sua griglia.} \\
\\ 
\\ 
\texttt{Qual è il tuo prossimo comando?} \\
            }
        }
    }
    & \\ \\

    \theutterance \stepcounter{utterance}  
    & & & \multicolumn{4}{p{0.6\linewidth}}{
        \cellcolor[rgb]{0.9,0.9,0.9}{
            \makecell[{{p{\linewidth}}}]{
                \texttt{\tiny{[GM$\langle$P2]}}
                \texttt{SPOSTA: L, (3, 4)} \\
            }
        }
    }
     \\ \\

    \theutterance \stepcounter{utterance}  
    & & & \multicolumn{2}{p{0.3\linewidth}}{
        \cellcolor[rgb]{0.9,0.9,0.9}{
            \makecell[{{p{\linewidth}}}]{
                \texttt{\tiny{[GM$|$GM]}}
                \texttt{SPOSTA: L, (3, 4)} \\
            }
        }
    }
    & & \\ \\

    \theutterance \stepcounter{utterance}  
    & & & \multicolumn{2}{p{0.3\linewidth}}{
        \cellcolor[rgb]{0.9,0.9,0.9}{
            \makecell[{{p{\linewidth}}}]{
                \texttt{\tiny{[GM$|$GM]}}
                \texttt{Penalità: (3,4) non è vuoto, ma contiene '┬'.} \\
\\ 
\texttt{Avete accumulato collettivamente 2 di 8 penalità.} \\
\texttt{Riprova!} \\
            }
        }
    }
    & & \\ \\

    \theutterance \stepcounter{utterance}  
    & & \multicolumn{4}{p{0.6\linewidth}}{
        \cellcolor[rgb]{0.9,0.9,0.9}{
            \makecell[{{p{\linewidth}}}]{
                \texttt{\tiny{[GM$\rangle$P2]}}
                \texttt{Penalità: (3,4) non è vuoto, ma contiene '┬'.} \\
\\ 
\texttt{Avete accumulato collettivamente 2 di 8 penalità.} \\
\texttt{Riprova!} \\
            }
        }
    }
    & \\ \\

    \theutterance \stepcounter{utterance}  
    & & & \multicolumn{4}{p{0.6\linewidth}}{
        \cellcolor[rgb]{0.9,0.9,0.9}{
            \makecell[{{p{\linewidth}}}]{
                \texttt{\tiny{[GM$\langle$P2]}}
                \texttt{COMUNICA: Mi scuso per l'errore: (3,4) non è vuoto sulla mia griglia — contiene '┬', una linea di connessione. Il mio piano di spostare L in (3,4) non è possibile come previsto. Devo rivedere la configurazione obiettivo.  } \\
\texttt{Per favore, conferma se puoi spostare i tuoi oggetti in alternative sicure, ad esempio:  } \\
\texttt{{-} P → (4,5)  } \\
\texttt{{-} L → (3,5)  } \\
\texttt{{-} C → (5,5)  } \\
\texttt{Queste posizioni sono tutte vuote sulla mia griglia (controllato: (4,5)=◌, (3,5)=◌, (5,5)=◌).  } \\
\texttt{Attendo la tua conferma prima di proseguire con gli spostamenti.} \\
            }
        }
    }
     \\ \\

    \theutterance \stepcounter{utterance}  
    & & & \multicolumn{2}{p{0.3\linewidth}}{
        \cellcolor[rgb]{0.9,0.9,0.9}{
            \makecell[{{p{\linewidth}}}]{
                \texttt{\tiny{[GM$|$GM]}}
                \texttt{COMUNICA: Mi scuso per l'errore: (3,4) non è vuoto sulla mia griglia — contiene '┬', una linea di connessione. Il mio piano di spostare L in (3,4) non è possibile come previsto. Devo rivedere la configurazione obiettivo.  } \\
\texttt{Per favore, conferma se puoi spostare i tuoi oggetti in alternative sicure, ad esempio:  } \\
\texttt{{-} P → (4,5)  } \\
\texttt{{-} L → (3,5)  } \\
\texttt{{-} C → (5,5)  } \\
\texttt{Queste posizioni sono tutte vuote sulla mia griglia (controllato: (4,5)=◌, (3,5)=◌, (5,5)=◌).  } \\
\texttt{Attendo la tua conferma prima di proseguire con gli spostamenti.} \\
            }
        }
    }
    & & \\ \\

    \theutterance \stepcounter{utterance}  
    & & & \multicolumn{2}{p{0.3\linewidth}}{
        \cellcolor[rgb]{0.9,0.9,0.9}{
            \makecell[{{p{\linewidth}}}]{
                \texttt{\tiny{[GM$|$GM]}}
                \texttt{Invalid format: tail is not empty: '} \\
\texttt{Per favore, conferma se puoi spostare i tuoi oggetti in alternative sicure, ad esempio:  } \\
\texttt{{-} P → (4,5)  } \\
\texttt{{-} L → (3,5)  } \\
\texttt{{-} C → (5,5)  } \\
\texttt{Queste posizioni sono tutte vuote sulla mia griglia (controllato: (4,5)=◌, (3,5)=◌, (5,5)=◌).  } \\
\texttt{Attendo la tua conferma prima di proseguire con gli spostamenti.} \\
            }
        }
    }
    & & \\ \\

    \theutterance \stepcounter{utterance}  
    & & \multicolumn{4}{p{0.6\linewidth}}{
        \cellcolor[rgb]{0.9,0.9,0.9}{
            \makecell[{{p{\linewidth}}}]{
                \texttt{\tiny{[GM$\rangle$P2]}}
                \texttt{Penalità: Il tuo messaggio non deve contenere nulla dopo il comando!} \\
\\ 
\texttt{Assicurati che la tua risposta contenga solo `COMUNICA: <MESSAGE>` o `SPOSTA: <OBJECT>, (<X>, <Y>)`, e nient'altro!} \\
\\ 
\texttt{Avete accumulato collettivamente 3 di 8 penalità.} \\
\texttt{Riprova!} \\
            }
        }
    }
    & \\ \\

    \theutterance \stepcounter{utterance}  
    & & & \multicolumn{4}{p{0.6\linewidth}}{
        \cellcolor[rgb]{0.9,0.9,0.9}{
            \makecell[{{p{\linewidth}}}]{
                \texttt{\tiny{[GM$\langle$P2]}}
                \texttt{COMUNICA: Mi scuso per l'errore. Propongo una nuova configurazione obiettivo: P in (4,5), L in (3,5), C in (5,5). Tutte queste celle sono vuote sulla mia griglia. Confermi?} \\
            }
        }
    }
     \\ \\

    \theutterance \stepcounter{utterance}  
    & & & \multicolumn{2}{p{0.3\linewidth}}{
        \cellcolor[rgb]{0.9,0.9,0.9}{
            \makecell[{{p{\linewidth}}}]{
                \texttt{\tiny{[GM$|$GM]}}
                \texttt{COMUNICA: Mi scuso per l'errore. Propongo una nuova configurazione obiettivo: P in (4,5), L in (3,5), C in (5,5). Tutte queste celle sono vuote sulla mia griglia. Confermi?} \\
            }
        }
    }
    & & \\ \\

    \theutterance \stepcounter{utterance}  
    & & \multicolumn{4}{p{0.6\linewidth}}{
        \cellcolor[rgb]{0.9,0.9,0.9}{
            \makecell[{{p{\linewidth}}}]{
                \texttt{\tiny{[P1$\langle$GM]}}
                \texttt{Spostato 'L' in (2,5) con successo. La tua griglia aggiornata appare così:} \\
\texttt{```} \\
\texttt{ 1234567} \\
\texttt{╔═╤═╤═╤═╗} \\
\texttt{║◌│◌│◌│◌║ 1} \\
\texttt{║◌└─┼─┴─╢ 2} \\
\texttt{║◌◌◌│◌◌◌║ 3} \\
\texttt{╟──┬┴───╢ 4} \\
\texttt{║◌L│P◌◌◌║ 5} \\
\texttt{║◌C│◌◌◌◌║ 6} \\
\texttt{║◌◌│◌◌◌◌║ 7} \\
\texttt{╚══╧════╝} \\
\texttt{```} \\
\texttt{Stai attualmente giocando il round 2 di massimo 12.} \\
\texttt{Avete accumulato collettivamente 3 di 8 penalità.} \\
\texttt{L'altro giocatore ha inviato questo messaggio:} \\
\\ 
\\ 
\texttt{"Mi scuso per l'errore. Propongo una nuova configurazione obiettivo: P in (4,5), L in (3,5), C in (5,5). Tutte queste celle sono vuote sulla mia griglia. Confermi?"} \\
\\ 
\\ 
\texttt{Qual è il tuo prossimo comando?} \\
            }
        }
    }
    & \\ \\

    \theutterance \stepcounter{utterance}  
    & \multicolumn{4}{p{0.6\linewidth}}{
        \cellcolor[rgb]{0.9,0.9,0.9}{
            \makecell[{{p{\linewidth}}}]{
                \texttt{\tiny{[P1$\rangle$GM]}}
                \texttt{COMUNICA: Confermo! La tua nuova configurazione (P in (4,5), L in (3,5), C in (5,5)) è perfetta. Sulla mia griglia, (4,5) è vuoto, (3,5) è vuoto, e (5,5) è già occupato da 'P' — quindi devo spostare il mio 'P' da (5,5) a (4,5), e poi il mio 'C' da (2,6) a (5,5). Inizierò con il movimento di 'P' verso (4,5).} \\
            }
        }
    }
    & & \\ \\

    \theutterance \stepcounter{utterance}  
    & & & \multicolumn{2}{p{0.3\linewidth}}{
        \cellcolor[rgb]{0.9,0.9,0.9}{
            \makecell[{{p{\linewidth}}}]{
                \texttt{\tiny{[GM$|$GM]}}
                \texttt{COMUNICA: Confermo! La tua nuova configurazione (P in (4,5), L in (3,5), C in (5,5)) è perfetta. Sulla mia griglia, (4,5) è vuoto, (3,5) è vuoto, e (5,5) è già occupato da 'P' — quindi devo spostare il mio 'P' da (5,5) a (4,5), e poi il mio 'C' da (2,6) a (5,5). Inizierò con il movimento di 'P' verso (4,5).} \\
            }
        }
    }
    & & \\ \\

    \theutterance \stepcounter{utterance}  
    & & \multicolumn{4}{p{0.6\linewidth}}{
        \cellcolor[rgb]{0.9,0.9,0.9}{
            \makecell[{{p{\linewidth}}}]{
                \texttt{\tiny{[GM$\rangle$P2]}}
                \texttt{Il tuo messaggio è stato inoltrato all'altro giocatore.} \\
\\ 
\\ 
\texttt{Stai attualmente giocando il round 3 di massimo 12.} \\
\texttt{Avete accumulato collettivamente 3 di 8 penalità.} \\
\texttt{L'altro giocatore ha inviato questo messaggio:} \\
\\ 
\\ 
\texttt{"Confermo! La tua nuova configurazione (P in (4,5), L in (3,5), C in (5,5)) è perfetta. Sulla mia griglia, (4,5) è vuoto, (3,5) è vuoto, e (5,5) è già occupato da 'P' — quindi devo spostare il mio 'P' da (5,5) a (4,5), e poi il mio 'C' da (2,6) a (5,5). Inizierò con il movimento di 'P' verso (4,5)."} \\
\\ 
\\ 
\texttt{Qual è il tuo prossimo comando?} \\
            }
        }
    }
    & \\ \\

    \theutterance \stepcounter{utterance}  
    & & & \multicolumn{4}{p{0.6\linewidth}}{
        \cellcolor[rgb]{0.9,0.9,0.9}{
            \makecell[{{p{\linewidth}}}]{
                \texttt{\tiny{[GM$\langle$P2]}}
                \texttt{SPOSTA: L, (3, 5)} \\
            }
        }
    }
     \\ \\

    \theutterance \stepcounter{utterance}  
    & & & \multicolumn{2}{p{0.3\linewidth}}{
        \cellcolor[rgb]{0.9,0.9,0.9}{
            \makecell[{{p{\linewidth}}}]{
                \texttt{\tiny{[GM$|$GM]}}
                \texttt{SPOSTA: L, (3, 5)} \\
            }
        }
    }
    & & \\ \\

    \theutterance \stepcounter{utterance}  
    & & & \multicolumn{2}{p{0.3\linewidth}}{
        \cellcolor[rgb]{0.9,0.9,0.9}{
            \makecell[{{p{\linewidth}}}]{
                \texttt{\tiny{[GM$|$GM]}}
                \texttt{Penalità: (3,5) non è vuoto, ma contiene '│'.} \\
\\ 
\texttt{Avete accumulato collettivamente 4 di 8 penalità.} \\
\texttt{Riprova!} \\
            }
        }
    }
    & & \\ \\

    \theutterance \stepcounter{utterance}  
    & & \multicolumn{4}{p{0.6\linewidth}}{
        \cellcolor[rgb]{0.9,0.9,0.9}{
            \makecell[{{p{\linewidth}}}]{
                \texttt{\tiny{[GM$\rangle$P2]}}
                \texttt{Penalità: (3,5) non è vuoto, ma contiene '│'.} \\
\\ 
\texttt{Avete accumulato collettivamente 4 di 8 penalità.} \\
\texttt{Riprova!} \\
            }
        }
    }
    & \\ \\

    \theutterance \stepcounter{utterance}  
    & & & \multicolumn{4}{p{0.6\linewidth}}{
        \cellcolor[rgb]{0.9,0.9,0.9}{
            \makecell[{{p{\linewidth}}}]{
                \texttt{\tiny{[GM$\langle$P2]}}
                \texttt{COMUNICA: (3,5) non è vuoto sulla mia griglia — contiene '│'. La mia configurazione di obiettivo deve essere aggiornata. Propongo: P in (4,5), L in (2,5), C in (5,5). Verifico: (4,5)=◌, (2,5)=◌, (5,5)=◌. Confermi questa nuova configurazione?} \\
            }
        }
    }
     \\ \\

    \theutterance \stepcounter{utterance}  
    & & & \multicolumn{2}{p{0.3\linewidth}}{
        \cellcolor[rgb]{0.9,0.9,0.9}{
            \makecell[{{p{\linewidth}}}]{
                \texttt{\tiny{[GM$|$GM]}}
                \texttt{COMUNICA: (3,5) non è vuoto sulla mia griglia — contiene '│'. La mia configurazione di obiettivo deve essere aggiornata. Propongo: P in (4,5), L in (2,5), C in (5,5). Verifico: (4,5)=◌, (2,5)=◌, (5,5)=◌. Confermi questa nuova configurazione?} \\
            }
        }
    }
    & & \\ \\

    \theutterance \stepcounter{utterance}  
    & & \multicolumn{4}{p{0.6\linewidth}}{
        \cellcolor[rgb]{0.9,0.9,0.9}{
            \makecell[{{p{\linewidth}}}]{
                \texttt{\tiny{[P1$\langle$GM]}}
                \texttt{Il tuo messaggio è stato inoltrato all'altro giocatore.} \\
\\ 
\\ 
\texttt{Stai attualmente giocando il round 3 di massimo 12.} \\
\texttt{Avete accumulato collettivamente 4 di 8 penalità.} \\
\texttt{L'altro giocatore ha inviato questo messaggio:} \\
\\ 
\\ 
\texttt{"(3,5) non è vuoto sulla mia griglia — contiene '│'. La mia configurazione di obiettivo deve essere aggiornata. Propongo: P in (4,5), L in (2,5), C in (5,5). Verifico: (4,5)=◌, (2,5)=◌, (5,5)=◌. Confermi questa nuova configurazione?"} \\
\\ 
\\ 
\texttt{Qual è il tuo prossimo comando?} \\
            }
        }
    }
    & \\ \\

    \theutterance \stepcounter{utterance}  
    & \multicolumn{4}{p{0.6\linewidth}}{
        \cellcolor[rgb]{0.9,0.9,0.9}{
            \makecell[{{p{\linewidth}}}]{
                \texttt{\tiny{[P1$\rangle$GM]}}
                \texttt{COMUNICA: Confermo la nuova configurazione: P in (4,5), L in (2,5), C in (5,5).  } \\
\texttt{Sulla mia griglia:  } \\
\texttt{{-} (4,5) è vuoto → posso spostarci 'P'  } \\
\texttt{{-} (2,5) è già occupato dal mio 'L' → perfetto, non muoverlo  } \\
\texttt{{-} (5,5) è occupato da 'P' → devo spostare 'P' da (5,5) a (4,5), poi 'C' da (2,6) a (5,5)  } \\
\\ 
\texttt{Inizierò subito spostando 'P' da (5,5) a (4,5).} \\
            }
        }
    }
    & & \\ \\

    \theutterance \stepcounter{utterance}  
    & & & \multicolumn{2}{p{0.3\linewidth}}{
        \cellcolor[rgb]{0.9,0.9,0.9}{
            \makecell[{{p{\linewidth}}}]{
                \texttt{\tiny{[GM$|$GM]}}
                \texttt{COMUNICA: Confermo la nuova configurazione: P in (4,5), L in (2,5), C in (5,5).  } \\
\texttt{Sulla mia griglia:  } \\
\texttt{{-} (4,5) è vuoto → posso spostarci 'P'  } \\
\texttt{{-} (2,5) è già occupato dal mio 'L' → perfetto, non muoverlo  } \\
\texttt{{-} (5,5) è occupato da 'P' → devo spostare 'P' da (5,5) a (4,5), poi 'C' da (2,6) a (5,5)  } \\
\\ 
\texttt{Inizierò subito spostando 'P' da (5,5) a (4,5).} \\
            }
        }
    }
    & & \\ \\

    \theutterance \stepcounter{utterance}  
    & & & \multicolumn{2}{p{0.3\linewidth}}{
        \cellcolor[rgb]{0.9,0.9,0.9}{
            \makecell[{{p{\linewidth}}}]{
                \texttt{\tiny{[GM$|$GM]}}
                \texttt{Invalid format: tail is not empty: '} \\
\texttt{Sulla mia griglia:  } \\
\texttt{{-} (4,5) è vuoto → posso spostarci 'P'  } \\
\texttt{{-} (2,5) è già occupato dal mio 'L' → perfetto, non muoverlo  } \\
\texttt{{-} (5,5) è occupato da 'P' → devo spostare 'P' da (5,5) a (4,5), poi 'C' da (2,6) a (5,5)  } \\
\\ 
\texttt{Inizierò subito spostando 'P' da (5,5) a (4,5).} \\
            }
        }
    }
    & & \\ \\

    \theutterance \stepcounter{utterance}  
    & & \multicolumn{4}{p{0.6\linewidth}}{
        \cellcolor[rgb]{0.9,0.9,0.9}{
            \makecell[{{p{\linewidth}}}]{
                \texttt{\tiny{[P1$\langle$GM]}}
                \texttt{Penalità: Il tuo messaggio non deve contenere nulla dopo il comando!} \\
\\ 
\texttt{Assicurati che la tua risposta contenga solo `COMUNICA: <MESSAGE>` o `SPOSTA: <OBJECT>, (<X>, <Y>)`, e nient'altro!} \\
\\ 
\texttt{Avete accumulato collettivamente 5 di 8 penalità.} \\
\texttt{Riprova!} \\
            }
        }
    }
    & \\ \\

    \theutterance \stepcounter{utterance}  
    & \multicolumn{4}{p{0.6\linewidth}}{
        \cellcolor[rgb]{0.9,0.9,0.9}{
            \makecell[{{p{\linewidth}}}]{
                \texttt{\tiny{[P1$\rangle$GM]}}
                \texttt{SPOSTA: P, (4,5)} \\
            }
        }
    }
    & & \\ \\

    \theutterance \stepcounter{utterance}  
    & & & \multicolumn{2}{p{0.3\linewidth}}{
        \cellcolor[rgb]{0.9,0.9,0.9}{
            \makecell[{{p{\linewidth}}}]{
                \texttt{\tiny{[GM$|$GM]}}
                \texttt{SPOSTA: P, (4,5)} \\
            }
        }
    }
    & & \\ \\

    \theutterance \stepcounter{utterance}  
    & & & \multicolumn{2}{p{0.3\linewidth}}{
        \cellcolor[rgb]{0.9,0.9,0.9}{
            \makecell[{{p{\linewidth}}}]{
                \texttt{\tiny{[GM$|$GM]}}
                \texttt{Penalità: (4,5) non è vuoto, ma contiene 'P'.} \\
\\ 
\texttt{Avete accumulato collettivamente 6 di 8 penalità.} \\
\texttt{Riprova!} \\
            }
        }
    }
    & & \\ \\

    \theutterance \stepcounter{utterance}  
    & & \multicolumn{4}{p{0.6\linewidth}}{
        \cellcolor[rgb]{0.9,0.9,0.9}{
            \makecell[{{p{\linewidth}}}]{
                \texttt{\tiny{[P1$\langle$GM]}}
                \texttt{Penalità: (4,5) non è vuoto, ma contiene 'P'.} \\
\\ 
\texttt{Avete accumulato collettivamente 6 di 8 penalità.} \\
\texttt{Riprova!} \\
            }
        }
    }
    & \\ \\

    \theutterance \stepcounter{utterance}  
    & \multicolumn{4}{p{0.6\linewidth}}{
        \cellcolor[rgb]{0.9,0.9,0.9}{
            \makecell[{{p{\linewidth}}}]{
                \texttt{\tiny{[P1$\rangle$GM]}}
                \texttt{SPOSTA: C, (5,5)} \\
            }
        }
    }
    & & \\ \\

    \theutterance \stepcounter{utterance}  
    & & & \multicolumn{2}{p{0.3\linewidth}}{
        \cellcolor[rgb]{0.9,0.9,0.9}{
            \makecell[{{p{\linewidth}}}]{
                \texttt{\tiny{[GM$|$GM]}}
                \texttt{SPOSTA: C, (5,5)} \\
            }
        }
    }
    & & \\ \\

    \theutterance \stepcounter{utterance}  
    & & & \multicolumn{2}{p{0.3\linewidth}}{
        \cellcolor[rgb]{0.9,0.9,0.9}{
            \makecell[{{p{\linewidth}}}]{
                \texttt{\tiny{[GM$|$GM]}}
                \texttt{Spostato 'C' in (5,5) con successo. La tua griglia aggiornata appare così:} \\
\texttt{```} \\
\texttt{ 1234567} \\
\texttt{╔═╤═╤═╤═╗} \\
\texttt{║◌│◌│◌│◌║ 1} \\
\texttt{║◌└─┼─┴─╢ 2} \\
\texttt{║◌◌◌│◌◌◌║ 3} \\
\texttt{╟──┬┴───╢ 4} \\
\texttt{║◌L│PC◌◌║ 5} \\
\texttt{║◌◌│◌◌◌◌║ 6} \\
\texttt{║◌◌│◌◌◌◌║ 7} \\
\texttt{╚══╧════╝} \\
\texttt{```} \\
            }
        }
    }
    & & \\ \\

    \theutterance \stepcounter{utterance}  
    & & \multicolumn{4}{p{0.6\linewidth}}{
        \cellcolor[rgb]{0.9,0.9,0.9}{
            \makecell[{{p{\linewidth}}}]{
                \texttt{\tiny{[GM$\rangle$P2]}}
                \texttt{Il tuo messaggio è stato inoltrato all'altro giocatore.} \\
\\ 
\\ 
\texttt{Stai attualmente giocando il round 4 di massimo 12.} \\
\texttt{Avete accumulato collettivamente 6 di 8 penalità.} \\
\texttt{L'altro giocatore ha spostato un oggetto sulla sua griglia.} \\
\\ 
\\ 
\texttt{Qual è il tuo prossimo comando?} \\
            }
        }
    }
    & \\ \\

    \theutterance \stepcounter{utterance}  
    & & & \multicolumn{4}{p{0.6\linewidth}}{
        \cellcolor[rgb]{0.9,0.9,0.9}{
            \makecell[{{p{\linewidth}}}]{
                \texttt{\tiny{[GM$\langle$P2]}}
                \texttt{SPOSTA: P, (4, 5)} \\
            }
        }
    }
     \\ \\

    \theutterance \stepcounter{utterance}  
    & & & \multicolumn{2}{p{0.3\linewidth}}{
        \cellcolor[rgb]{0.9,0.9,0.9}{
            \makecell[{{p{\linewidth}}}]{
                \texttt{\tiny{[GM$|$GM]}}
                \texttt{SPOSTA: P, (4, 5)} \\
            }
        }
    }
    & & \\ \\

    \theutterance \stepcounter{utterance}  
    & & & \multicolumn{2}{p{0.3\linewidth}}{
        \cellcolor[rgb]{0.9,0.9,0.9}{
            \makecell[{{p{\linewidth}}}]{
                \texttt{\tiny{[GM$|$GM]}}
                \texttt{Spostato 'P' in (4,5) con successo. La tua griglia aggiornata appare così:} \\
\texttt{```} \\
\texttt{ 1234567} \\
\texttt{╔═╤═╤═╤═╗} \\
\texttt{║◌│◌│L│◌║ 1} \\
\texttt{║◌└─┼─┴─╢ 2} \\
\texttt{║◌◌◌│◌◌C║ 3} \\
\texttt{╟──┬┴───╢ 4} \\
\texttt{║◌◌│P◌◌◌║ 5} \\
\texttt{║◌◌│◌◌◌◌║ 6} \\
\texttt{║◌◌│◌◌◌◌║ 7} \\
\texttt{╚══╧════╝} \\
\texttt{```} \\
            }
        }
    }
    & & \\ \\

    \theutterance \stepcounter{utterance}  
    & & \multicolumn{4}{p{0.6\linewidth}}{
        \cellcolor[rgb]{0.9,0.9,0.9}{
            \makecell[{{p{\linewidth}}}]{
                \texttt{\tiny{[P1$\langle$GM]}}
                \texttt{Spostato 'C' in (5,5) con successo. La tua griglia aggiornata appare così:} \\
\texttt{```} \\
\texttt{ 1234567} \\
\texttt{╔═╤═╤═╤═╗} \\
\texttt{║◌│◌│◌│◌║ 1} \\
\texttt{║◌└─┼─┴─╢ 2} \\
\texttt{║◌◌◌│◌◌◌║ 3} \\
\texttt{╟──┬┴───╢ 4} \\
\texttt{║◌L│PC◌◌║ 5} \\
\texttt{║◌◌│◌◌◌◌║ 6} \\
\texttt{║◌◌│◌◌◌◌║ 7} \\
\texttt{╚══╧════╝} \\
\texttt{```} \\
\texttt{Stai attualmente giocando il round 4 di massimo 12.} \\
\texttt{Avete accumulato collettivamente 6 di 8 penalità.} \\
\texttt{L'altro giocatore ha spostato un oggetto sulla sua griglia.} \\
\\ 
\\ 
\texttt{Qual è il tuo prossimo comando?} \\
            }
        }
    }
    & \\ \\

    \theutterance \stepcounter{utterance}  
    & \multicolumn{4}{p{0.6\linewidth}}{
        \cellcolor[rgb]{0.9,0.9,0.9}{
            \makecell[{{p{\linewidth}}}]{
                \texttt{\tiny{[P1$\rangle$GM]}}
                \texttt{SPOSTA: P, (4,5)} \\
            }
        }
    }
    & & \\ \\

    \theutterance \stepcounter{utterance}  
    & & & \multicolumn{2}{p{0.3\linewidth}}{
        \cellcolor[rgb]{0.9,0.9,0.9}{
            \makecell[{{p{\linewidth}}}]{
                \texttt{\tiny{[GM$|$GM]}}
                \texttt{SPOSTA: P, (4,5)} \\
            }
        }
    }
    & & \\ \\

    \theutterance \stepcounter{utterance}  
    & & & \multicolumn{2}{p{0.3\linewidth}}{
        \cellcolor[rgb]{0.9,0.9,0.9}{
            \makecell[{{p{\linewidth}}}]{
                \texttt{\tiny{[GM$|$GM]}}
                \texttt{Penalità: (4,5) non è vuoto, ma contiene 'P'.} \\
\\ 
\texttt{Avete accumulato collettivamente 7 di 8 penalità.} \\
\texttt{Riprova!} \\
            }
        }
    }
    & & \\ \\

    \theutterance \stepcounter{utterance}  
    & & \multicolumn{4}{p{0.6\linewidth}}{
        \cellcolor[rgb]{0.9,0.9,0.9}{
            \makecell[{{p{\linewidth}}}]{
                \texttt{\tiny{[P1$\langle$GM]}}
                \texttt{Penalità: (4,5) non è vuoto, ma contiene 'P'.} \\
\\ 
\texttt{Avete accumulato collettivamente 7 di 8 penalità.} \\
\texttt{Riprova!} \\
            }
        }
    }
    & \\ \\

    \theutterance \stepcounter{utterance}  
    & \multicolumn{4}{p{0.6\linewidth}}{
        \cellcolor[rgb]{0.9,0.9,0.9}{
            \makecell[{{p{\linewidth}}}]{
                \texttt{\tiny{[P1$\rangle$GM]}}
                \texttt{SPOSTA: L, (3,5)} \\
            }
        }
    }
    & & \\ \\

    \theutterance \stepcounter{utterance}  
    & & & \multicolumn{2}{p{0.3\linewidth}}{
        \cellcolor[rgb]{0.9,0.9,0.9}{
            \makecell[{{p{\linewidth}}}]{
                \texttt{\tiny{[GM$|$GM]}}
                \texttt{SPOSTA: L, (3,5)} \\
            }
        }
    }
    & & \\ \\

    \theutterance \stepcounter{utterance}  
    & & & \multicolumn{2}{p{0.3\linewidth}}{
        \cellcolor[rgb]{0.9,0.9,0.9}{
            \makecell[{{p{\linewidth}}}]{
                \texttt{\tiny{[GM$|$GM]}}
                \texttt{Penalità: (3,5) non è vuoto, ma contiene '│'.} \\
\\ 
\texttt{Avete accumulato collettivamente 8 di 8 penalità.} \\
\texttt{Riprova!} \\
            }
        }
    }
    & & \\ \\

    \theutterance \stepcounter{utterance}  
    & & \multicolumn{4}{p{0.6\linewidth}}{
        \cellcolor[rgb]{0.9,0.9,0.9}{
            \makecell[{{p{\linewidth}}}]{
                \texttt{\tiny{[P1$\langle$GM]}}
                \texttt{Penalità: (3,5) non è vuoto, ma contiene '│'.} \\
\\ 
\texttt{Avete accumulato collettivamente 8 di 8 penalità.} \\
\texttt{Riprova!} \\
            }
        }
    }
    & \\ \\

    \theutterance \stepcounter{utterance}  
    & \multicolumn{4}{p{0.6\linewidth}}{
        \cellcolor[rgb]{0.9,0.9,0.9}{
            \makecell[{{p{\linewidth}}}]{
                \texttt{\tiny{[P1$\rangle$GM]}}
                \texttt{COMUNICA: finito?} \\
            }
        }
    }
    & & \\ \\

    \theutterance \stepcounter{utterance}  
    & & & \multicolumn{2}{p{0.3\linewidth}}{
        \cellcolor[rgb]{0.9,0.9,0.9}{
            \makecell[{{p{\linewidth}}}]{
                \texttt{\tiny{[GM$|$GM]}}
                \texttt{COMUNICA: finito?} \\
            }
        }
    }
    & & \\ \\

    \theutterance \stepcounter{utterance}  
    & & \multicolumn{4}{p{0.6\linewidth}}{
        \cellcolor[rgb]{0.9,0.9,0.9}{
            \makecell[{{p{\linewidth}}}]{
                \texttt{\tiny{[GM$\rangle$P2]}}
                \texttt{Spostato 'P' in (4,5) con successo. La tua griglia aggiornata appare così:} \\
\texttt{```} \\
\texttt{ 1234567} \\
\texttt{╔═╤═╤═╤═╗} \\
\texttt{║◌│◌│L│◌║ 1} \\
\texttt{║◌└─┼─┴─╢ 2} \\
\texttt{║◌◌◌│◌◌C║ 3} \\
\texttt{╟──┬┴───╢ 4} \\
\texttt{║◌◌│P◌◌◌║ 5} \\
\texttt{║◌◌│◌◌◌◌║ 6} \\
\texttt{║◌◌│◌◌◌◌║ 7} \\
\texttt{╚══╧════╝} \\
\texttt{```} \\
\texttt{Stai attualmente giocando il round 5 di massimo 12.} \\
\texttt{Avete accumulato collettivamente 8 di 8 penalità.} \\
\texttt{L'altro giocatore ha inviato questo messaggio:} \\
\\ 
\\ 
\texttt{"finito?"} \\
\\ 
\\ 
\texttt{Qual è il tuo prossimo comando?} \\
            }
        }
    }
    & \\ \\

    \theutterance \stepcounter{utterance}  
    & & & \multicolumn{4}{p{0.6\linewidth}}{
        \cellcolor[rgb]{0.9,0.9,0.9}{
            \makecell[{{p{\linewidth}}}]{
                \texttt{\tiny{[GM$\langle$P2]}}
                \texttt{COMUNICA: finito!} \\
            }
        }
    }
     \\ \\

    \theutterance \stepcounter{utterance}  
    & & & \multicolumn{2}{p{0.3\linewidth}}{
        \cellcolor[rgb]{0.9,0.9,0.9}{
            \makecell[{{p{\linewidth}}}]{
                \texttt{\tiny{[GM$|$GM]}}
                \texttt{COMUNICA: finito!} \\
            }
        }
    }
    & & \\ \\

    \theutterance \stepcounter{utterance}  
    & & & \multicolumn{2}{p{0.3\linewidth}}{
        \cellcolor[rgb]{0.9,0.9,0.9}{
            \makecell[{{p{\linewidth}}}]{
                \texttt{\tiny{[GM$|$GM]}}
                \texttt{true} \\
            }
        }
    }
    & & \\ \\

    \theutterance \stepcounter{utterance}  
    & & & \multicolumn{2}{p{0.3\linewidth}}{
        \cellcolor[rgb]{0.9,0.9,0.9}{
            \makecell[{{p{\linewidth}}}]{
                \texttt{\tiny{[GM$|$GM]}}
                \texttt{* success: True} \\
\texttt{* lose: False} \\
\texttt{* aborted: False} \\
\texttt{{-}{-}{-}{-}{-}{-}{-}} \\
\texttt{* Shifts: 2.00} \\
\texttt{* Max Shifts: 4.00} \\
\texttt{* Min Shifts: 2.00} \\
\texttt{* End Distance Sum: 7.83} \\
\texttt{* Init Distance Sum: 13.65} \\
\texttt{* Expected Distance Sum: 12.57} \\
\texttt{* Penalties: 8.00} \\
\texttt{* Max Penalties: 8.00} \\
\texttt{* Rounds: 5.00} \\
\texttt{* Max Rounds: 12.00} \\
\texttt{* Object Count: 3.00} \\
            }
        }
    }
    & & \\ \\

    \theutterance \stepcounter{utterance}  
    & & & \multicolumn{2}{p{0.3\linewidth}}{
        \cellcolor[rgb]{0.9,0.9,0.9}{
            \makecell[{{p{\linewidth}}}]{
                \texttt{\tiny{[GM$|$GM]}}
                \texttt{Initial Grids:} \\
\texttt{```} \\
\texttt{Player 1:     Player 2:} \\
\texttt{ 1234567       1234567} \\
\texttt{╔═╤═╤═╤═╗     ╔═╤═╤═╤═╗} \\
\texttt{║◌│◌│◌│◌║ 1   ║◌│◌│L│◌║ 1} \\
\texttt{║◌└─┼─┴─╢ 2   ║◌└─┼─┴─╢ 2} \\
\texttt{║◌◌◌│◌◌◌║ 3   ║◌◌◌│◌◌C║ 3} \\
\texttt{╟──┬┴───╢ 4   ╟──┬┴───╢ 4} \\
\texttt{║◌◌│P◌◌◌║ 5   ║◌◌│◌◌◌◌║ 5} \\
\texttt{║LC│◌◌◌◌║ 6   ║◌◌│◌P◌◌║ 6} \\
\texttt{║◌◌│◌◌◌◌║ 7   ║◌◌│◌◌◌◌║ 7} \\
\texttt{╚══╧════╝     ╚══╧════╝} \\
\texttt{```} \\
\\ 
\texttt{Final Grids:} \\
\texttt{```} \\
\texttt{Player 1:     Player 2:} \\
\texttt{ 1234567       1234567} \\
\texttt{╔═╤═╤═╤═╗     ╔═╤═╤═╤═╗} \\
\texttt{║◌│◌│◌│◌║ 1   ║◌│◌│L│◌║ 1} \\
\texttt{║◌└─┼─┴─╢ 2   ║◌└─┼─┴─╢ 2} \\
\texttt{║◌◌◌│◌◌◌║ 3   ║◌◌◌│◌◌C║ 3} \\
\texttt{╟──┬┴───╢ 4   ╟──┬┴───╢ 4} \\
\texttt{║◌L│PC◌◌║ 5   ║◌◌│P◌◌◌║ 5} \\
\texttt{║◌◌│◌◌◌◌║ 6   ║◌◌│◌◌◌◌║ 6} \\
\texttt{║◌◌│◌◌◌◌║ 7   ║◌◌│◌◌◌◌║ 7} \\
\texttt{╚══╧════╝     ╚══╧════╝} \\
\texttt{```} \\
            }
        }
    }
    & & \\ \\

\end{supertabular}
}

\end{document}
